The verified statements below separate geometric variation from an
abstract lift through the boundary map, and separate a vanishing theorem
from the exact information supplied by its proposed formal argument.
For cote 125, the canonical module is denoted by $\omega_X$; it is
not silently replaced by the top exterior power of Kahler differentials.
The original pages have been compared with the supplied transcriptions
at every doubtful reading used below \cite{C-ms122,C-ms125}.

\subsection{Conventions and the equivariant boundary triangle}
We use cohomological complexes: $H^i(K[a])=H^{i+a}(K)$.
For a commutative coefficient ring $\Lambda$ we put
$D_\Lambda K=R\operatorname{Hom}_\Lambda(K,\Lambda)$.
For $\pi=\langle t\rangle\cong\mathbb Z$, duals carry the contragredient action.
A ``perfect'' complex means perfect over $\Lambda$, unless another ring is named.
The unqualified tensor and dual symbols below refer to derived operations when complexes occur.

\begin{saresult}{C.001}{Proved}{Factorization and its exact ambiguity}
Let $R\xrightarrow{u}K\xrightarrow{\rho}L\xrightarrow{\delta}R[1]$
be a distinguished triangle in $D(\Lambda[t,t^{-1}])$, and suppose that
$t-1=0$ as an endomorphism of $L$ in that category.
Then there is a morphism $V:K\to R$ with $uV=t-1$.
For any one such $V_0$, the entire set of choices is
\[
 V_0+\operatorname{im}\bigl(
 \operatorname{Hom}(K,L[-1])\xrightarrow{\delta[-1]\circ-}
 \operatorname{Hom}(K,R)\bigr).
\]
In particular, existence does not imply uniqueness.
This applies to the triangle of cote 122, scan 022, with
$R=D_\Lambda K[-2n]$, once its boundary action is trivial as an object of the derived category.
\end{saresult}
\begin{proof}
Equivariance gives $\rho(t-1)=(t-1)\rho=0$.
Apply $\operatorname{Hom}(K,-)$ to the triangle. Exactness of
\[
 \operatorname{Hom}(K,L[-1])\longrightarrow
 \operatorname{Hom}(K,R)\xrightarrow{u\circ-}
 \operatorname{Hom}(K,K)\xrightarrow{\rho\circ-}
 \operatorname{Hom}(K,L)
\]
gives both assertions. The displayed image describes differences of lifts, not merely an upper bound on their ambiguity. We have used triviality in the derived category; triviality of each cohomology representation alone is a weaker assertion over a general coefficient ring.
\end{proof}

% The following results are inserted here for incremental assembly. Their IDs,
% rather than their position in this working file, determine the final order.

\begin{saresult}{C.002}{Refuted}{The uniqueness assertion of Lead L4}
The hypotheses of C.001, including perfectness, Poincare--Lefschetz duality,
and a trivial boundary action, do not force $V$ to be unique.
There is a geometric counterexample in every integer $n\geq2$:
\[
 M=S^n\times D^n,\qquad \partial M=S^n\times S^{n-1},
 \qquad h=\operatorname{id}_M,
\]
with coefficients in any field $\Lambda$.
The actual variation of this identity monodromy is zero, whereas the factorization equation admits nonzero solutions.
\end{saresult}
\begin{proof}
Give both factors their standard orientations. The finite CW pair gives perfect
cochain complexes $K=R\Gamma(M,\Lambda)$ and
$R=R\Gamma(M,\partial M;\Lambda)$; the action is literally trivial on the complexes.
Kunneth and the relative cohomology of the disk give
\[
 H^*(K)=\Lambda\text{ in degrees }0,n,
 \qquad H^*(R)=\Lambda\text{ in degrees }n,2n.
\]
The relative-to-absolute map $u$ is zero in cohomology: its degree-$n$
class is pulled back from $H^n(D^n,S^{n-1})$ and dies in
$H^n(D^n)=0$, while the target in degree $2n$ is zero.
Over a field every complex splits, as a complex of vector spaces, into
its cohomology and an acyclic summand. All actions here are trivial,
so the splitting and its contracting homotopies are equivariant.
Consequently $u=0$ in $D(\Lambda[t,t^{-1}])$.
The composite
\[
 K\longrightarrow H^n(K)[-n]=\Lambda[-n]
 \xrightarrow{1}H^n(R)[-n]=\Lambda[-n]\longrightarrow R
\]
is an equivariant derived morphism $V_1$ inducing the identity in degree $n$;
therefore it is nonzero. Both $0$ and $V_1$ satisfy $uV=t-1=0$.
The boundary triangle and its self-duality are those of this actual compact
oriented manifold, so those additional hypotheses have all been verified.
This refutes the universal uniqueness claim, not the existence of a distinguished
variation coming from specified geometric monodromy.
\end{proof}

\begin{saresult}{C.003}{Proved}{The geometric variation is a specified lift}
Let $(M,B)$ be a compact triangulable manifold with boundary $B$, and let
$h:M\to M$ be a homeomorphism that is the identity on $B$.
Let $t$ act on singular chains by $h_\#$ and on cochains by pullback
$h^*$ (this choice fixes the generator convention).
There is an equivariant cochain map
\[
 V:C^*(M;\Lambda)\longrightarrow C^*(M,B;\Lambda),
 \qquad (Vc)([z])=c(h_\#z-z).
\]
If $u:C^*(M,B)\to C^*(M)$ is the usual inclusion, then
$uV=h^*-1$ and $Vu=h^*_{\mathrm{rel}}-1$ as actual cochain maps.
The induced derived map agrees with classical variation in SGA~7~II,
XIV, 3.1.9--3.1.10, after matching its transport-of-structure convention.
\end{saresult}
\begin{proof}
If $z$ is changed by a chain in $B$, its difference $h_\#z-z$ is unchanged
because $h$ fixes every simplex in $B$. Hence the formula is well defined
on relative chains. The chain map $h_\#-1$ commutes with the boundary,
so precomposition gives a cochain map. It commutes with $h^*$ because
$h_\#$ commutes with $h_\#-1$. The two identities are immediate on every
chain and thus hold before passing to a derived category.
Replacing $h$ by an isotopy relative to $B$ changes $h_\#-1$ by the prism
chain homotopy relative to $B$; the induced variation is therefore unchanged
in the derived category. The construction in XIV, 3.1.9 uses an equivariant
resolution and factors $T-1$ into the kernel of restriction to $B$.
Singular cochains provide such a resolution on this triangulable pair,
and the map just constructed is exactly that kernel factorization.
Proposition 3.1.10 identifies this construction with the variation from
nearby and vanishing cycles. This comparison is not a uniqueness argument
inside a bare triangulated category.

The source locators used above are \cite[XIV, 3.1.9--3.1.10, pp.~142--145]{C-sga7ii}.
\end{proof}

\begin{saresult}{C.007}{Proved}{Boundary transpose signs and the Seifert form}
Orient a compact smooth $2n$-manifold $M$ and its boundary $B$ by the
outward-normal-first convention. Write $r:H^i(M)\to H^i(B)$ for
restriction and $\delta:H^{2n-i-1}(B)\to H^{2n-i}(M,B)$ for the
connecting map. With ordinary cohomology in the first slot of each
duality pairing,
\[
 \langle b,\delta a\rangle_M
       =(-1)^i\langle rb,a\rangle_B.
 \tag{C.7.1}
\]
Thus the degree-$n$ arrows $\rho$ and $\rho'$ in scan 022 are transposes
with sign $(-1)^n$. For the whole graded triangle the sign is $(-1)^i$,
not one scalar independent of degree. Using the relative class in the
first slot instead gives the sign-free identity
$\langle\delta a,b\rangle_M=\langle a,rb\rangle_B$.

For an isolated hypersurface Milnor page let $H=H_n(M;\mathbb Z)$,
$Q=H_n(M,B;\mathbb Z)$, $v:Q\to H$ be homological variation for
positive monodromy, and $\epsilon=(-1)^n$. The pairing $(q,x)$ between
$Q$ and $H$ is the oriented intersection number. In the Seifert convention
\[
 L(x,y)=-\epsilon\,(v^{-1}y,x),
 \tag{C.7.2}
\]
the intersection pairing $I$ represented by the middle map $u$ satisfies
\[
 I=L+\epsilon L^{\mathsf t},\qquad
 L(x,Ty)=-\epsilon L(y,x).
 \tag{C.7.3}
\]
In column-vector bases dual for $(q,x)$, these are
$V=-\epsilon L^{-1}$ and $T=-\epsilon L^{-1}L^{\mathsf t}$.
These conventions specify the signs completely.
\end{saresult}
\begin{proof}
First take differential-form representatives. Extend a closed form $a$
from the boundary to a collar and then to a form $\widetilde a$ on $M$;
the relative form $d\widetilde a$ represents $\delta a$. If $b$ is closed,
Stokes' formula gives
\[
 \int_M d\widetilde a\wedge b=\int_B a\wedge rb,
 \qquad
 \int_M b\wedge d\widetilde a=(-1)^i\int_B rb\wedge a.
\]
This proves (C.7.1) over characteristic zero. The same proof on singular
cochains uses the identity
$d(b\smile\widetilde a)=(-1)^i b\smile d\widetilde a$ and evaluation
on a relative fundamental chain whose boundary is the boundary
fundamental chain. It therefore proves the identity over $\mathbb Z$
and with all coefficient rings, including characteristic two. A collar
permits the extension cochain; changing it changes the connecting
representative by a relative coboundary.

For the Seifert identities the precise classical source is Lamotke,
\emph{Die Homologie isolierter Singularit\"aten}, \S3.2,
(3.2.5)--(3.2.6), \S3.3--3.4, and \S6.1,
(6.1.3)--(6.1.6), pp.~30--33 and 40--41.
Its isomorphism theorem applies to the open book with total sphere and
gives that $v$ is an integral isomorphism in degree $n$.
The two identities from \S3.2, in the present notation, are
\[
 I(va,vb)=-(a,vb)-\epsilon(b,va),\qquad
 (T_Qb,va)=-\epsilon(a,vb).
 \tag{C.7.4}
\]
Here $T_Q$ is relative monodromy and $vT_Q=Tv$.
Substitute $x=va$, $y=vb$ in (C.7.2).
The first equation of (C.7.4) is exactly
$I(x,y)=L(x,y)+\epsilon L(y,x)$.
The second gives
$L(x,Ty)=-\epsilon(T_Qb,va)=(a,vb)=-\epsilon L(y,x)$.
Thus the cited identities imply the two formulas with no unrecorded
transpose or monodromy inversion. If $q_i,e_i$ are dual bases,
$L(e_i,e_j)=-\epsilon(V^{-1})_{ij}$, proving the matrix formula.
Since the matrix of $j:H\to Q$ is $I^{\mathsf t}=\epsilon I$,
the identity $vj=T-1$ becomes
$(-\epsilon L^{-1})(\epsilon I)=T-1$, which is precisely (C.7.3).
Dualizing gives $uV_{mathrm{coh}}=T^{\mathsf t}-1$ on cohomology.
Using transport monodromy $(h^{-1})^*$ instead of $h^*$ amounts to
replacing the chosen geometric generator by its inverse throughout.

The source locators used above are \cite[3.1--3.4 and 6.1--6.2, pp.~29--33, 40--41]{C-lamotke1975}.
\end{proof}

\begin{saresult}{C.008}{Known}{The tame etale triangle, its variation, and every Tate twist}
Let $S$ be an excellent strictly henselian trait of residue characteristic
$p$, let $f:X\to S$ be separated of finite type, and let $x$ be a geometric
closed point of the special fibre. Suppose that near $x$ the generic
fibre is smooth and pure of dimension $n$, and that $f$ is smooth on the
special fibre away from $x$. Take $\Lambda=\mathbb Z/\ell^a$ with
$\ell\ne p$, or a finite field of characteristic $\ell$.
Let $I$ be inertia, $P_w$ its wild pro-$p$ subgroup, and $I_t=I/P_w$.
On the strict localization at $x$ put
\[
 K=(R\Psi\Lambda)_x,\quad
 R=R\Gamma_{\{x\}}(X_s,R\Psi\Lambda),\quad
 L=R\Gamma(X_s\setminus\{x\},\Lambda).
\]
There is an $I$-equivariant triangle
\[
 R\xrightarrow{u}K\longrightarrow L^{\rm triv}\longrightarrow R[1],
 \qquad R\simeq D_\Lambda K(-n)[-2n].
 \tag{C.8.1}
\]
Consequently $K_t=K^{P_w}$ has the triangle in the continuous derived
category of $\Lambda[I_t]$-modules
\[
 D_\Lambda K_t(-n)[-2n]\xrightarrow{u_t}K_t
 \longrightarrow L^{\rm triv}\longrightarrow.
 \tag{C.8.2}
\]
For a chosen $\sigma\in I_t$, canonical nearby/vanishing-cycle variation
gives $V_\sigma:K_t\to D_\Lambda K_t(-n)[-2n]$ with
$u_tV_\sigma=\sigma-1$. This choice is geometric; it is not the unique
abstract lift furnished by the triangle. If wild inertia acts trivially,
$K_t=K$ and (C.8.2) is the requested tame version for the whole fibre.

For a bounded perfect tame complex, continuous group duality gives
\[
 R\Gamma(I_t,D_\Lambda K_t)
 \simeq D_\Lambda R\Gamma(I_t,K_t)(-1)[-1].
 \tag{C.8.3}
\]
Hence the left term after derived tame invariants in (C.8.2) is
$D_\Lambda R\Gamma(I_t,K_t)(-n-1)[-2n-1]$.
For trivial $L$ the right term is
$L\oplus L(-1)[-1]$. The shifts and twists $(-n)$ and $(-n-1)$
are therefore forced, rather than interchangeable conventions.
\end{saresult}
\begin{proof}[Derivation from the opened etale theorems]
Localization for the closed immersion of $x$ gives the first triangle.
Smooth base change identifies $R\Psi\Lambda$ with the constant sheaf
off $x$, compatibly with the trivial inertia action; this is SGA~7~II,
XIII, 2.1.4--2.1.5. The same paragraph identifies $K$ with the cohomology
of the geometric generic fibre of the strict localization.
Gabber's nearby-cycle duality theorem, in Illusie,
\emph{Autour du theoreme de monodromie locale}, Theorem 4.2 and
(4.3.6), pp.~45--46, applies to this separated finite-type morphism.
It says $D_{X_s}R\Psi\Lambda\simeq R\Psi D_{X_{\bar\eta}}\Lambda$.
Smooth purity on the generic fibre makes the right side
$R\Psi\Lambda(n)[2n]$.
The identity $i^!D_{X_s}=D_\Lambda i^*$, applied to the immersion
$i:x\hookrightarrow X_s$, gives exactly the second formula of (C.8.1).
Constructibility and finite Tor dimension in the same theorem make the
stalk complexes perfect over the stated coefficient rings.

Taking $P_w$-invariants is exact: every continuous finite quotient of
$P_w$ has order a power of $p$, and its averaging idempotent is defined
over $\Lambda$. The idempotent is self-adjoint for the contragredient
pairing, so duality commutes with its image. Passage to invariants in
(C.8.1) therefore proves (C.8.2). This is also the tame-cycle construction
of SGA~7~I, I, 2.7.1--2.7.3.
SGA~7~II, XIII, 1.4.2--1.4.3 constructs variation before forgetting the
specialization complex. At an isolated point, XIII, (2.4.5.1) gives
\[
 (R\Phi\Lambda)_x\longrightarrow
 R\Gamma_{\{x\}}(R\Psi\Lambda).
\]
Its composite with $K\to(R\Phi\Lambda)_x$ has composite $\sigma-1$
with $u$, by (2.4.6.1). Exact wild invariants give the assertion for
$I_t$. There is no dependence on a chosen lift of $\sigma$ to
$I$: after taking wild invariants, the variation cocycle on
$P_w$ is a continuous additive homomorphism into the finite
$\ell$-primary group of maps between the resulting perfect
complexes. Such a homomorphism from a pro-$p$ group is zero.
The cocycle identity of XIII, 1.4.3.4, therefore gives the same
variation for two lifts. Equivalently one performs the same
mapping-cone construction on the tame nearby-cycle complex;
the tame action is abelian, so the construction is equivariant
for the whole $I_t$.

Finally $I_t\cong\prod_{q\ne p}\mathbb Z_q(1)$.
First average the prime-to-$\ell$ factor. A topological generator of the
remaining $\mathbb Z_\ell(1)$ represents its cohomology by the two-term
complex $[K_t\xrightarrow{\sigma-1}K_t]$ in group degrees $0,1$,
with the degree-one term carrying the dual orientation $\Lambda(-1)$.
Dualizing reverses those degrees; shifting by $[-1]$ restores them.
The dual action is $(\sigma^{-1})^{\mathsf t}$, and the equality
\[
 (\sigma^{-1})^{\mathsf t}-1
   =-(\sigma^{-1})^{\mathsf t}(\sigma^{\mathsf t}-1)
\]
identifies the two differentials. Totalizing with a perfect representative
of $K_t$ proves (C.8.3). For a trivial action the differential is zero.
This also explains the extra shift in scan 024 and why etale degree one
adds a Tate twist. Changing the generator changes the coordinate used
on $\mathbb Z_\ell(1)$, not the intrinsically twisted formula.

The source locators used above are \cite[I, 2.7.1--2.7.3]{C-sga7i}; \cite[XIII, 1.4 and 2.4]{C-sga7ii}; \cite[Theorem~4.2 and (4.3.6), pp.~45--46]{C-illusie1994}.
\end{proof}

\begin{saresult}{C.009}{Refuted}{Tame inertia cannot silently replace full inertia}
Smooth total space and an isolated non-smooth point do not in themselves
make nearby cycles tame. Consequently the assertion that the whole
$K$ carries a tame-inertia action requires a tameness hypothesis;
without it, (C.8.2) applies to $K^{P_w}$.
\end{saresult}
\begin{proof}
Let $k$ be algebraically closed of characteristic $p>0$ and consider
\[
 f:\operatorname{Spec}k[x,(1+x)^{-1}]\longrightarrow\mathbb A^1_k,
 \qquad t=\frac{x^p}{1+x}.
\]
The source is a smooth curve. Its derivative is
$-x^p/(1+x)^2$, so its only non-smooth point is $x=0$.
The equation $x^p-tx-t=0$ shows that the map is finite of degree $p$;
$1+x$ is automatically a unit in this finite algebra because the
equation takes a nonzero constant value at $x=-1$.
Over the generic point the derivative with respect to $x$ is $-t$,
so the extension is separable. At $t=0$ there is one point with
ramification degree $p$. After strict henselization the geometric
generic fibre has $p$ points and inertia acts transitively on them.
If wild inertia acted trivially, this transitive permutation action
would factor through a finite quotient of tame inertia, whose order is
prime to $p$. Such a group cannot have an orbit of length $p$.
Thus wild inertia acts nontrivially already on $H^0(K)$.
This example has fibre dimension zero; it disproves the automatic
tameness assertion in general without making an unsupported claim
about a higher-dimensional suspension. C.008 supplies the version
valid in every dimension.
\end{proof}

\begin{saresult}{C.004}{Proved under stated hypotheses}{Closing the missing vanishing for $n\geq2$}
Let $M$ be a connected compact oriented $2n$-manifold with nonempty boundary
$B$, where $n\geq2$. Suppose that an open book with page $M$, binding $B$,
and monodromy fixed on $B$ has total space $W=S^{2n+1}$. Assume that
$B$ is $(n-2)$-connected and that $H^j(M;\Lambda)=0$ for $j>n$, where
$\Lambda$ is a field. Then $H^{n-1}(M;\Lambda)=0$.
These hypotheses hold for the compact Milnor page of an isolated critical
point of a holomorphic function on a smooth complex $(n+1)$-fold.
\end{saresult}
\begin{proof}
Write $P=W\setminus B$, which fibres over $S^1$ with fibre homotopy
equivalent to $M$. Poincare--Lefschetz duality and the assumed upper
vanishing give $H^{n-1}(M,B)=H^{n+1}(M)^\vee=0$.
Thus the boundary exact sequence gives an injection
\[
 H^{n-1}(M)\hookrightarrow H^{n-1}(B).
\]
It is monodromy equivariant; the action on its target is the identity.
Hence $T=1$ on $H^{n-1}(M)$.

For $n\geq3$, the Gysin sequence for the oriented codimension-two embedding
$B\subset S^{2n+1}$ gives
\[
 H^{n-1}(P)\cong H^{n-2}(B)=0.
\]
The Wang exact sequence gives a surjection
$H^{n-1}(P)\twoheadrightarrow
\ker(T-1:H^{n-1}(M)\to H^{n-1}(M))$.
The latter kernel is the whole group, which therefore vanishes.

For $n=2$, the Gysin sequence gives $H^1(P)\cong H^0(B)=\Lambda$.
Let $p:P\to S^1$ be the open-book projection. A positive normal meridian
of $B$ maps to $S^1$ with degree one by the definition of an open book.
Consequently the residue of $p^*$ applied to the positive generator of
$H^1(S^1;\Lambda)$ equals $1\in H^0(B;\Lambda)$.
It follows that $p^*:H^1(S^1;\Lambda)\to H^1(P;\Lambda)$ is an isomorphism.
The degree-one Wang sequence is
\[
 0\longrightarrow H^1(S^1;H^0(M))
 \xrightarrow{p^*}H^1(P)
 \longrightarrow H^1(M)^T\longrightarrow0.
\]
Connectedness of $M$ identifies its left group with $\Lambda$.
Thus $H^1(M)^T=0$, and $T=1$ on $H^1(M)$ proves the claim.

For a holomorphic isolated critical point, Milnor's Theorem 4.8 supplies
the fibration, Lemma 6.1 the compact page and boundary, Theorem 5.1 the
upper vanishing, and Theorem 5.2 the connectivity of $B$.
The smooth total-space link is a sphere. Connectedness of $M$ when
$n\geq2$ also follows from the boundary sequence:
$H^0(M,B)=H^{2n}(M)^\vee=0$ and
$H^1(M,B)=H^{2n-1}(M)^\vee=0$, whence
$H^0(M)\cong H^0(B)=\Lambda$.
The positive meridian condition is the local model $f/|f|$ near the binding.
No middle-dimensional vanishing theorem for the fibre has been used as an input.

The source locators used above are \cite[Theorems~4.8, 5.1--5.2, Lemma~6.1, Lemma~8.4]{C-milnor1968}.
\end{proof}

\begin{saresult}{C.005}{Refuted}{The unreduced assertion in dimension one}
The statement $H^{n-1}(M;\Lambda)=0$ in C.1(d) and Lead L5 is false if
$n=1$ is included literally. The correct uniform assertion is
$\widetilde H^{n-1}(M;\Lambda)=0$.
\end{saresult}
\begin{proof}
Take $f:\mathbb C^2\to\mathbb C$, $f(x,y)=xy$.
Its derivative $(y,x)$ vanishes only at the origin, and the total space
$\mathbb C^2$ is smooth. For $0<|\delta|<\epsilon^2/2$, its compact Milnor fibre
is parameterized by $x\mapsto(x,\delta/x)$ with
\[
 |x|^2+|\delta|^2/|x|^2\leq\epsilon^2.
\]
Equivalently $a\leq |x|^2\leq b$, where
$a,b=(\epsilon^2\mp\sqrt{\epsilon^4-4|\delta|^2})/2$ are positive and $a<b$.
The fibre is an annulus. Thus $H^0(M;\Lambda)=\Lambda\ne0$, while
$\widetilde H^0(M;\Lambda)=0$.
For an arbitrary isolated plane-curve critical point, connectedness of
the Milnor page follows without using its middle Betti number:
the open-book complement of one closed page fibres over an open interval
and is homotopy equivalent to the page; Alexander duality identifies
its reduced $H_0$ with $\widetilde H^{2}(M)$, which is zero by the
one-dimensional CW bound. This proves the reduced assertion for $n=1$.
Together with C.004 it gives the stated uniform correction.

The source locators used above are \cite[Corollaries~6.2--6.3, pp.~56--57]{C-milnor1968}.
\end{proof}

\begin{saresult}{C.006}{Proved}{The exact scope of the missing vanishing argument}
For an isolated complex hypersurface critical point with smooth total
space and complex fibre dimension $n\geq1$,
$\widetilde H^j(M;A)=0$ for $j<n$ and every abelian coefficient group $A$.
In particular the unreduced assertion requested in C.1(d) holds when
$n\geq2$, and must be reduced when $n=1$.
The boundary duality triangle, its six-term middle sequence, the upper
cohomological bound, and $(n-2)$-connectedness of the boundary, taken by
themselves, do not imply $H^{n-1}(M)=0$.
\end{saresult}
\begin{proof}
For a field, the triangle and Poincare--Lefschetz duality give
$H^j(M)\simeq H^j(\partial M)$ for $j\leq n-2$:
the two adjacent relative groups are dual to groups of degrees strictly
greater than $n$, which vanish by Milnor's upper-dimensional bound.
The link connectivity therefore gives reduced vanishing in that range.
C.004 gives the remaining degree $n-1$ when $n\geq2$, and C.005 gives
the reduced degree-zero statement when $n=1$.
The page has finite CW type. Apply the field conclusion to $\mathbb Q$
and to $\mathbb F_p$ for every prime $p$. If a finitely generated reduced
integral homology group in a degree below $n$ were nonzero, its free part
would be detected over $\mathbb Q$, or its primary torsion would be
detected over the corresponding $\mathbb F_p$, by the homology universal
coefficient sequence. Thus those integral groups vanish. The cohomology
universal coefficient sequence then proves the assertion for every $A$.

For the insufficiency assertion take, for any $n\geq2$,
\[
 M=S^{n-1}\times D^{n+1},\qquad
 B=\partial M=S^{n-1}\times S^n,\qquad h=\operatorname{id}.
\]
This is a compact oriented $2n$-manifold with boundary, its boundary is
$(n-2)$-connected, and it retracts onto an $(n-1)$-sphere. It has the
genuine relative cochain triangle, equivariantly with trivial action,
and Poincare--Lefschetz duality makes that triangle self-dual.
Nevertheless $H^{n-1}(M;\Lambda)=\Lambda$ for every nonzero field
$\Lambda$. Its six-term middle sequence is exactly
\[
0\longrightarrow\Lambda\xrightarrow{1}\Lambda\longrightarrow0
 \longrightarrow0\longrightarrow\Lambda\xrightarrow{1}\Lambda
 \longrightarrow0.
\]
The associated open book is the boundary of $M\times D^2$, namely
$S^{n-1}\times S^{n+2}$ after rounding corners. It is not a sphere.
This verifies every formal input just listed and isolates the missing
geometric input: the sphere total space and its meridian/Gysin sequence.
It does not furnish, and is not claimed to furnish, an isolated complex
singularity with smooth total space. For that geometric problem C.004
uses precisely the additional input and completes the argument.
\end{proof}

\begin{saresult}{C.017}{Refuted}{Five literal errors and their precise corrections in cote 122}
The following assertions, if read literally in scans 003, 007, 022,
030, 032 and 034, are false. Each has the correction specified below.
\begin{enumerate}
\item For a merely paracompact manifold with boundary, compactly
supported interior cohomology cannot replace relative cohomology in
the ordinary boundary triangle. On a compact finite-type manifold it
can. Without finite-type hypotheses, duality must also retain its
correct dual direction rather than identify a vector space with its
double dual.
\item In scan 007 the complementary manifold is
$Q=W\setminus\mathring P$, not the written $P\setminus\mathring P$.
\item In scan 022 an equivariant map
$D_\Lambda K[-2n]\to K$ is classified by
$H^{2n}(R\Gamma(\pi,K\otimes^L K))$, not in general by
$H^{2n}(K\otimes^L K)$.
\item For a Milnor open-book complement $P=W\setminus B$, the upper
bound is $H^i(P)=0$ for $i>n+1$. The stronger bound $i>n$ in one
reading of scan 030 fails. The upper bound $H^i(M)=0$ for $i>n$
for the fibre is true independently.
\item The arrow $H^{n-1}(B)\to H^n(M)^\vee$ in scan 032 need not
be surjective. In scan 034 the list of possible nonzero degrees of
$B=V_s$ must include its top degree $2n-1$.
\end{enumerate}
\end{saresult}
\begin{proof}
(1) The insertion ``para'' is visible in the archive image 003.
For $P=\mathbb R$ with empty boundary,
$H_c^0(\mathbb R;\Lambda)=0$ and $H^0(\mathbb R;\Lambda)=\Lambda$.
Thus the displayed boundary sequence cannot be exact with compact
supports. The exact replacement is
\[
 R\Gamma(P,j_!\Lambda_{\mathring P})\longrightarrow
 R\Gamma(P,\Lambda)\longrightarrow R\Gamma(\partial P,\Lambda).
\]
The first term is relative cohomology, or cohomology with supports in
the interior which are closed in $P$. Such supports are compact when
$P$ is compact. For finite coefficients over a field, Verdier/Poincare
duality naturally identifies ordinary cohomology with the dual of
compactly supported cohomology in complementary degree. Reversing
this identification requires finite-dimensionality. For example,
on a countable discrete zero-dimensional manifold and over
$\mathbb F_2$, compactly supported $H^0$ is a countable-dimensional
direct sum, whereas $\operatorname{Hom}_{\mathbb F_2}
(H^0,\mathbb F_2)$ is the dual of an infinite product and is much
larger. Compact triangulable pages have finite cohomology and avoid
both problems.

(2) The ink on scan 007 itself writes $Q=P-\mathring P$.
Take $W=S^1$ and $P$ a proper closed semicircle. That written $Q$
consists of two endpoints, so $P\cup Q=P\ne W$.
With $Q=W\setminus\mathring P$ and $P$ closed, codimension zero,
and locally flat at its boundary, collars identify the two sides of
the boundary with $B\times[0,\epsilon)$ and
$B\times(-\epsilon,0]$. This proves
$W=P\cup Q$, $P\cap Q=\partial P=\partial Q$, and both manifold
assertions. The resulting localization sequence is Mayer--Vietoris
after replacing the two closed pieces by collar enlargements.

(3) Perfect tensor--Hom adjunction in the equivariant derived category
gives exactly the group written in the corrected assertion. Forgetting
equivariance gives its edge map to ordinary tensor cohomology, which
need not be injective. For an explicit example take a field $k$,
$\pi=\mathbb Z$, $n\geq2$, and the trivial-action complex
$K=k\oplus k[-(2n-1)]$. Then
$H^{2n}(K\otimes K)=0$, but
\[
 H^{2n}(R\Gamma(\pi,K\otimes K))
 \supset H^1(\pi,H^{2n-1}(K\otimes K))=k^2.
\]
The two-term resolution of the infinite cyclic group proves the
displayed equality and exhibits nonzero equivariant maps whose
underlying maps vanish.

(4) The Wang sequence for a page of CW dimension at most $n$
gives the bound $n+1$ for its mapping torus. The critical point
$f(x,y)=xy$ has annular page and identity action on $H^1$;
its sphere complement is the complement of the Hopf link, or directly
the mapping torus in the open book. The Wang sequence gives
$H^2(P;\Lambda)=\Lambda\ne0$, although $n=1$.
Thus the reading $i>n$ is false. If the retraced bound is instead
read $i>n+1$, it is correct, but passage to a limit of those identical
inequalities does not on its own yield the sharper degree $n$.
Milnor's Theorem 5.1 supplies that sharper fibre bound.

(5) Take $f(z_0,z_1,z_2)=z_0^2+z_1^2+z_2^2$.
Its only critical point is the origin and its total space is smooth.
Writing $z=x+iy$ on the real positive fibre gives
$|x|^2-|y|^2=\delta$, $x\cdot y=0$.
The map $(x,y)\mapsto(x/|x|,y)$ identifies the compact fibre with
a disk tangent bundle of $S^2$, with its radius rescaled by the ball
cutoff. Its boundary is the unit tangent bundle $SO(3)=\mathbb RP^3$.
Over $\mathbb Q$, $H^1(B)=0$ and $H^2(M)^\vee=\mathbb Q$.
Thus the arrow labeled ``surj'' is not onto. Its middle intersection
map is multiplication by $\pm2$, the Euler number of $TS^2$.
Finally any connected closed oriented $(2n-1)$-manifold has top
cohomology $\Lambda$, so the omission in scan 034 is impossible.
For a genuine isolated singularity with $n=3$ take
$\sum_{i=0}^3z_i^2$: its link is such a closed oriented
five-manifold, and $H^5(B;\Lambda)=\Lambda$, outside the written
list $0,2,3$. The corrected list is $0,n-1,n,2n-1$.

The source locators used above are \cite[scans 003, 007, 022, 030, 032, 034]{C-ms122}.
\end{proof}

\begin{saresult}{C.010}{Proved}{The open-book package and the complete concordance for 122/003--035}
Let $M$ be a compact oriented $2n$-manifold with nonempty boundary $B$,
where $n\geq1$, let $h:M\to M$ preserve orientation and fix a collar of $B$, and set
\[
 P=(M\times[0,1])/((x,1)\sim(hx,0)),\quad
 Q=B\times D^2,\quad W=P\cup_{B\times S^1}Q.
\]
The gluing uses the identity on $B$ and the positive angular coordinate.
It gives a closed oriented $(2n+1)$-manifold. With coefficients in a
field, all the boundary, support, Mayer--Vietoris, duality and residue
statements on scans 003--034 hold with the corrections in C.017 and
the conventions below. The same derived triangles hold over any
commutative ring when the finite complexes and orientation twists are
retained. In particular, if $W$ is a cohomology sphere, then
\[
 H^i(P)\xrightarrow{\ \beta\ }H^{i-1}(B)
 \quad\hbox{is an isomorphism for }i\notin\{0,2n\},
 \tag{C.10.1}
\]
and
\[
 H^0(P)=\Lambda,\qquad
 0\longrightarrow H^{2n}(P)\longrightarrow H^{2n-1}(B)
 \xrightarrow{\int_B}\Lambda\longrightarrow0.
 \tag{C.10.2}
\]
Thus $H^{2n}(P)=0$ exactly when $B$ is connected. No unqualified
assertion about degree zero of a one-dimensional Milnor fibre is
implicit here.
\end{saresult}
\begin{proof}
Collars allow one to push $B$ a small positive distance into $M$;
the resulting map and the identity are joined by the collar
homotopy. This proves $\mathring M\hookrightarrow M$ is a homotopy
equivalence, also for locally constant coefficients. The same
argument in each half-ball shows the local-support sheaves
$\mathcal H^i_B(F)$ on $M$ vanish. The boundary triangle is
\[
 R\Gamma(M,B;F)\longrightarrow R\Gamma(M;F)
 \longrightarrow R\Gamma(B;F)\longrightarrow.
 \tag{C.10.3}
\]
For compact $M$ its first term is $R\Gamma_c(\mathring M;F)$.
Poincare--Lefschetz duality identifies it with
$D_\Lambda R\Gamma(M;F^\vee\otimes\operatorname{or}_M)[-2n]$,
when the coefficients are finite locally free and the complexes are
perfect. All restriction and connecting arrows are the actual
arrows of this triangle; the transposition sign is C.007.

Here is a precise version of the collar-excision argument on scan
005. Choose two nested closed cores $M_0\subset\mathring M_1\subset
\mathring M$ by deleting collar strips. All the inclusions of pairs
obtained by replacing a collar strip by a longer or shorter strip
are homotopy equivalences, by rescaling its interval coordinate.
Excision and these equivalences give
\[
 R\Gamma_{M_0}(M;F)\simeq R\Gamma(M,B;F)
 \simeq R\Gamma(M_1,\partial M_1;F).
\]
If the notation $M/B$ is used with constant coefficients, it means
\emph{reduced} quotient cohomology; with a general local system it
means the relative complex (C.10.3). Ordinary unreduced quotient
cohomology would already fail in degree zero for an interval.
The literal inclusion $(M_0,\partial M_0)\to(M,M\setminus M_0)$
is not a map of pairs: its boundary lies in $M_0$. One must use a
second core as above, or replace the complement by
$M\setminus\mathring M_0$ and then use the collar comparison.

For any compact codimension-zero collared decomposition
$W=P\cup Q$ with common boundary $R$, collar enlargement replaces
$P,Q$ by open sets retracting onto them, with intersection
retracting onto $R$. The ordinary Mayer--Vietoris triangle is
therefore
\[
 R\Gamma(W)\longrightarrow R\Gamma(P)\oplus R\Gamma(Q)
 \xrightarrow{\,r_P-r_Q\,}R\Gamma(R)\longrightarrow.
 \tag{C.10.4}
\]
The support triangle for $P\subset W$ is identified by excision
with
\[
 R\Gamma(P,\partial P)\longrightarrow R\Gamma(W)
 \longrightarrow R\Gamma(Q)\longrightarrow,
 \tag{C.10.5}
\]
and there is the analogous triangle with $P,Q$ exchanged.
Verdier duality, or the chain-level Poincare pairings, exchanges
these two triangles. It exchanges (C.10.4) with its compact-support
version. Applying the same constructions to the inclusions of
the common boundary produces the commutative localization
diagrams on scans 009, 011 and 017. These statements are identities
of the localization and excision maps, not merely a comparison of
their dimensions. The codimension-one Thom class of the common
boundary identifies its support triangle with (C.10.4); changing
the chosen normal direction changes both residue maps by the same
sign.

Now retain the open-book notation in the statement. Put
$K=R\Gamma(M)$, $R_M=R\Gamma(M,B)$ and $L=R\Gamma(B)$.
The monodromy is literally the identity on $B$, so (C.10.3) is the
triangle $R_M\to K\to L^{\rm triv}$ in
$D(\Lambda[\mathbb Z])$ used in C.001. A two-term free resolution
of the trivial cyclic-group module gives
\[
 R\Gamma(\mathbb Z,K)\simeq
 \operatorname{Cone}(T-1:K\to K)[-1]\simeq R\Gamma(P).
 \tag{C.10.6}
\]
One can see the second isomorphism directly by cutting the mapping
torus into two overlapping interval bundles: its Cech complex is
the same cone. The identical construction for the boundary gives
$R\Gamma(B\times S^1)=L\oplus L[-1]$ and for relative cochains
gives $R\Gamma(P,\partial P)$. This proves all three interpretations
of the triangles on scan 020. Its degree-one summand is a shift,
not a Tate twist, in ordinary topology.

Write restriction to $B\times S^1$ as
\[
 (\alpha,\beta):R\Gamma(P)\longrightarrow L\oplus L[-1].
\]
Thom excision in $Q=B\times D^2$ identifies the localization
triangle of $(W,W\setminus B)$ with
\[
 L[-2]\longrightarrow R\Gamma(W)\longrightarrow
 R\Gamma(P)\xrightarrow{\beta}L[-1].
 \tag{C.10.7}
\]
Indeed, on each normal disk the connecting map sends the positive
angular class to the relative fundamental class. Tensoring this
disk calculation with cochains on $B$ proves the claimed equality
with $\beta$, including its normalization. The other triangle is
\[
 R\Gamma(P,\partial P)\longrightarrow R\Gamma(W)
 \xrightarrow{\alpha\circ r_P}L\longrightarrow.
 \tag{C.10.8}
\]
Restriction from $(W,Q)$ to $(P,\partial P)$ is the excision
isomorphism in its first term. This proves the assertion about
$u_0,u_1,u_2$ on scan 022. The two triangles (C.10.7) and (C.10.8)
are dual under the orientation of $W$.

Dualizing the two-term cyclic resolution yields
\[
 R\Gamma(\mathbb Z,D_\Lambda K)
 \simeq D_\Lambda R\Gamma(\mathbb Z,K)[-1].
\]
The inverse-transpose identity for $T-1$ in the proof of C.008
checks the differential and the contragredient action. Consequently
the shift in scan 024 is $[-2n-1]$, as written. In the order
``class on $B$ first'', its two pairings are precisely
\[
 \lambda(a,x)=\int_B a\smile\beta(x),\qquad
 \mu(a,x)=\int_B a\smile\alpha(x).
 \tag{C.10.9}
\]
Their degrees are $-2n$ and $-(2n-1)$, respectively. Perfect
duality on $B$ recovers $\lambda'=\beta$ and $\mu'=\alpha$.
The other order introduces the usual Koszul sign. Thus scan 026
recovers the boundary restriction, rather than an additional
independent operation. Dual local coefficients and orientation
sheaves give exactly the variant stated there.

The long exact sequence of (C.10.7) gives (C.10.1). The complement
of a nonempty codimension-two collared submanifold in a connected
manifold is connected: perturb a path in its normal disks to avoid
their origins. Its top cohomology is zero. The top Gysin map is
integration over the components of $B$, so it is the summation map
$\Lambda^{\pi_0B}\to\Lambda$. This proves (C.10.2), including the
disconnected-link case in dimension one. Conversely, for a connected
closed oriented $W$, these residue isomorphisms and the injectivity
in degree $2n$ force $H^q(W)=0$ for $0<q<2n+1$, by exactness on
either side of each degree. This is the precise converse in the
margin of scan 028.

For an isolated hypersurface singularity in a smooth total space,
Milnor's fibration and compact-page results supply this open book.
His CW bound and link connectivity, followed by Poincare duality
on $B$, give the bounds and link-degree list on scans 028--032.
The long exact sequence of (C.10.3), using the CW bound, is exactly
the six-term sequence of scan 032, with no additional surjectivity
at its middle arrow. C.004--C.006 prove the missing middle-degree
vanishing and the edge cases. This completes the derivation of all
the topological statements used in the concordance below.

The source locators used above are \cite{C-ms122,C-milnor1968,C-sga7i,C-sga7ii,C-lamotke1975,C-illusie1994}.
\end{proof}

\paragraph{Statement-by-statement concordance.}
The following list distinguishes a source containing a statement
from a source supplying ingredients for the derivation just given.
``No separate counterpart'' means precisely that the page is a
formal topological construction or bookkeeping, not an assertion
that should be attributed to a numbered theorem on singularities.
The original SGA~7~I scan, the original SGA~7~II scan and the supplied
translation have been compared at the locators used here.

\begin{description}
\item[003.] Interior inclusion, vanishing of support on the boundary,
the relative triangle, and Poincare--Lefschetz pairing are proved in
C.010. Milnor, Lemma~6.1, supplies the compact collared page needed
in the application; SGA~7~II, XIV, 3.1.2(c),(d) and 3.1.7(B),
3.1.8.5--3.1.8.6, identify the corresponding compact-support
complex. Neither is a justification of the erroneous general
paracompact compact-support formulation; see C.017(1).

\item[005.] Core excision and the four collar comparisons are the
explicit collar argument in C.010. The support identification used
for nearby cycles is SGA~7~II, XIV, 3.1.8.5--3.1.8.6. There is no
separate Milnor theorem for the diagram. Relative, reduced-quotient
and unreduced-quotient readings have the distinction proved above.

\item[007.] Correct complementary piece, dimension, orientation
sheaf, codimension-one localization and its Mayer--Vietoris
interpretation are (C.10.4)--(C.10.5). The written complement is
false, C.017(2). These are formal topology; Milnor, Lemma~6.1,
provides the geometric boundary, and SGA~7~II, XIV, 3.1.2,
the boundary fibration, rather than this particular diagram.

\item[009 and 011.] Both localization sequences, their thickening
and excision identifications, the diamond, and the dual pairs
$\Phi,\Phi',\Psi,\Psi'$ are (C.10.4)--(C.10.5) and their duals.
The support comparison used by SGA is XIII, 2.4.2--2.4.6, and XIV,
3.1.7--3.1.10. No theorem of Milnor asserts this entire diamond;
his Lemma~6.1, Corollary~6.2 and Alexander-duality argument in Corollary~6.3
use the same compact-page geometry.

\item[013 and 015.] The general double-fibration construction is
valid when the two boundary identifications are given; collars
make the pushout a manifold and prove the interior inclusions are
homotopy equivalences. In the Milnor specialization the disk is
real two-dimensional, its boundary is $S^1$, and the other base is
$B=\partial M$. Milnor, Theorem~4.8, Theorem~5.11 and Lemma~6.1,
provide the sphere fibration, its tube comparison and the compact
page. SGA~7~II, XIV, 3.1.1--3.1.3 and 3.1.5, give the analogous
box, disk, boundary trivialization and monodromy representative.
Boundary triviality is available here; the interior homotopy
equivalence itself only needs a collar.

\item[017.] The geometric dictionary $W,P,Q,R;M,N,S,T$, the
restrictions and the displayed cohomology diamond specialize the
preceding two entries. They are proved by the same pair maps and
excision. The counts of possible diagram faces and sequences are
working bookkeeping, with no specified equivalence relation on
sequences and no counterpart in Milnor or SGA~7. They are not
used as a mathematical count in this treatment.

\item[020.] Triangle (a) is the equivariant boundary triangle;
(b) is its derived cyclic invariants, (c) its forgetful image, and
(d) the codimension-two Gysin triangle. Equations
(C.10.3),(C.10.6),(C.10.7) prove all four, including the equality
of the residue and $\beta$. Milnor's monodromy and Wang sequence
appear in Chapter~8, especially Lemma~8.4 and its sequence on p.~67.
SGA~7~I, I, 2.2--2.6, gives specialization and vanishing-cycle
sequences; SGA~7~II, XIII, 2.1.2--2.1.8 and 2.4.5--2.4.6, gives
their derived maps. These are analogous constructions, not an
identification of a punctured trait with an ordinary circle.

\item[022.] Triangle (e), the three comparison arrows and the
excision isomorphism are (C.10.8). Poincare--Lefschetz gives $u$;
the equivariant classifying group requires the correction
C.017(3). The transposition sign is C.007. SGA~7~II, XIII,
1.4.3.1--1.4.3.4, 2.1.8.7--2.1.8.8 and 2.4.5.1--2.4.6.1,
construct variation and factor monodromy; XIV, 3.1.9--3.1.10,
compares it with the chain map. SGA~7 does not say an abstract
factorization through $u$ is unique. Milnor's book supplies the
page and middle-dimensional homology, not this equivariant
uniqueness claim. The exact Seifert comparison is Lamotke,
3.2.5--3.2.6 and 6.1.3--6.1.6, with conventions in C.007.

\item[024.] Forgetful self-duality, duality after cyclic invariants,
the extra shift, and the identifications of $\lambda',\mu'$ are
(C.10.6)--(C.10.9). There is no additional Milnor theorem for these
derived formulas. The tame-inertia analogue uses SGA~7~I, I,
0.3.1 and 2.7.1--2.7.3; its duality and twists are derived in
C.008. General nearby-cycle duality in that derivation is
Gabber's theorem as stated by Illusie, Theorem~4.2 and (4.3.6),
pp.~45--46; attributing that full generality to SGA~7 would be
anachronistic.

\item[026.] The boundary pairing recovers restriction, (d) and (e)
are dual, and dual local coefficients give the same constructions:
these are (C.10.7)--(C.10.9) and C.007. The sphere Gysin
isomorphism follows from (C.10.7); compare Milnor's Alexander
duality in Corollary~6.3, pp.~56--57. SGA~7~II, XIV, 3.1.7--3.1.10, is the
nearby-cycle support counterpart, not a printed statement of this
whole self-dual diagram.

\item[028.] The top residue sequence, the disconnected-link
qualification, the converse sphere criterion, and the low/high
dual bounds are C.010 and C.004. Milnor, Theorem~5.2, gives
$(n-2)$-connectedness of the link, Theorem~5.1 the upper CW
bound; these imply the stated local Lefschetz range and the
mapping-torus bound $n+1$. SGA~7~I, I, Theorem~4.2, gives the
etale upper bound, and Theorem~4.5 and Corollaries~4.6--4.7
(characteristic zero as printed) give lower vanishing for a
relative complete intersection. Its proof explicitly invokes
SGA~2, XIV, 5.6, not XII--XIII. SGA~7~II, XIII, 2.1.13, p.~106, records
the upper bound again.

\item[030.] The sphere-complement degrees, the four nontrivial
pieces of the Gysin sequence, and the separate $n=1$ top sequence
are C.010. The fibre upper bound is Milnor, Theorem~5.1, and
SGA~7~I, I, Theorem~4.2; the displayed limit argument alone does
not improve a degree bound. C.017(4) treats both readings of the
retraced $n$ or $n+1$. The relative sequence at the foot of the
page is (C.10.3) with Poincare--Lefschetz duality.

\item[032.] Low and high restriction isomorphisms and the six-term
sequence are derived in C.010; the extra ``surj'' label is false
by C.017(5). Milnor, Corollaries~6.2--6.3, Lemma~6.4 and Theorem~6.5,
pp.~56--58, proves the
stronger reduced vanishing and bouquet theorem. C.004--C.006
derive the cohomological vanishing using only the explicitly named
weaker inputs. SGA~7~I, I, 4.5--4.7, supplies the characteristic-zero
complete-intersection etale result. SGA~7~II, XIV, 3.1.7(A),
identifies vanishing cycles with reduced Milnor cohomology;
XV, Proposition~3.1.2, is a more specific quadratic-singularity
statement and must not be cited as the general isolated case.

\item[034.] The renamed sphere-complement residue statements are
C.10.1--C.10.2. The final list for $V_s$ must contain $2n-1$;
the literal list is refuted in C.017(5). Milnor, Theorem~5.2,
Poincare duality on the compact link, and C.010 give the corrected
list. SGA~7~II, XV, 3.1.2, 3.2.1--3.2.3 and 3.3.5, instead give
the quadratic local cohomology, twists and Picard--Lefschetz
variation; their hypotheses are narrower than those of this page.

\item[035.] This is a separate unfinished note on finiteness of
inductive adic systems and stability under kernels or cokernels.
Neither the category denoted by the doubtful $\operatorname{ad}$
notation nor its proposed finiteness condition is defined on
these pages. The crossed-out example stops before defining $Q$.
It therefore has no identifiable proposition about a Milnor fibre,
and no counterpart in Milnor or in SGA~7~I or XIII--XV. No
stability assertion is inferred from this incomplete fragment.
The ordinary finite-adic completion theorem used in C.016 is a
different, explicitly stated assertion.
\end{description}
Scans 004, 006, 008, 010, 012, 014, 016, 018, 021, 023, 025,
027, 029, 031 and 033 are reused printed faces according to the
source; scan 019 has no mathematical text in the original image.
This accounts for every scan in the requested range.

The printed-page coordinates for the cited SGA passages are as follows:
SGA~7~I, Expose~I, 0.3.1 is on pp.~3--4, 2.2--2.6 on
pp.~8--9, 2.7.1--2.7.3 on p.~10, Theorem~4.2 on p.~14,
and Theorem~4.5 with Corollaries~4.6--4.7 on p.~15.
SGA~7~II, XIII, 1.4.3 is on p.~98, 2.1.2--2.1.8 on
pp.~99--103, 2.1.13 on p.~106, and 2.4.2--2.4.6 on
pp.~112--115; the two factorization formulas 2.4.5.1 and
2.4.6.1 are on p.~113. Expose~XIV, 3.1.1--3.1.3 is on
pp.~135--136, 3.1.5 on pp.~137--138, 3.1.7 on p.~139,
3.1.8 on pp.~140--142, 3.1.9 on pp.~142--143, and
Proposition~3.1.10 on pp.~143--145. Expose~XV,
Proposition~3.1.2 is on p.~188, 3.2.1--3.2.3 on
pp.~188--190, and Proposition~3.3.5 on p.~192.
Milnor's Theorem~4.8 is on p.~43, Theorems~5.1--5.2 on
p.~45, Theorem~5.11 on p.~53, and Lemma~6.1 on pp.~55--56.

\subsection{Cote 125: projective and punctured local duality}

\begin{saresult}{C.011}{Proved}{A finite full-support module detects nonzero tensors}
Let $A$ be a commutative noetherian ring, let $E$ be a finite $A$-module
with $\operatorname{Supp}_A(E)=\operatorname{Spec}(A)$, and let $N$ be any
$A$-module. Then $E\otimes_A N=0$ implies $N=0$.
\end{saresult}
\begin{proof}
First we show that $\operatorname{Hom}_A(E,Q)\ne0$ for every nonzero
$A$-module $Q$. A nonzero cyclic submodule of $Q$ is noetherian, hence
contains a submodule isomorphic to $A/\mathfrak p$ for some prime
$\mathfrak p$. Nakayama's lemma applied to $E_{\mathfrak p}$ shows that
$(E/\mathfrak pE)\otimes_{A/\mathfrak p}\operatorname{Frac}(A/\mathfrak p)$
is nonzero. Choose a nonzero functional on this finite-dimensional vector
space. Clearing the denominators of its values on finitely many generators
produces a nonzero map $E/\mathfrak pE\to A/\mathfrak p$.
Composition with $A/\mathfrak p\hookrightarrow Q$ proves the assertion.
Now take the character module
$Q=\operatorname{Hom}_{\mathbb Z}(N,\mathbb Q/\mathbb Z)$.
Since $\mathbb Q/\mathbb Z$ is an injective cogenerator of abelian groups,
$Q$ is nonzero if $N$ is nonzero. Tensor--Hom adjunction gives
\[
 \operatorname{Hom}_A(E,Q)
 =\operatorname{Hom}_{\mathbb Z}(E\otimes_A N,\mathbb Q/\mathbb Z).
\]
Thus a zero tensor product forces $Q=0$ and then $N=0$.
\end{proof}

\begin{saresult}{C.012}{Proved}{Negative twists of a full-support sheaf test cohomological dimension}
Let $X$ be a projective scheme over a field $k$, let $\mathcal O_X(1)$
be a very ample line bundle, let $Y\subseteq X$ be closed, and put $U=X\setminus Y$.
Let $E$ be a coherent sheaf with support all of $X$, and let $m\geq1$.
Then the following are equivalent:
\begin{enumerate}
\item $H^p(U,F)=0$ for every coherent sheaf $F$ on $U$ and every $p\geq m$;
\item $H^p(U,E(-N)|_U)=0$ for all $p\geq m$ and all sufficiently large integers $N$.
\end{enumerate}
There is no local-freeness assumption on $E$.
\end{saresult}
\begin{proof}
The forward implication is immediate. For the converse choose a standard
graded noetherian $k$-algebra $A$ with $X=\operatorname{Proj}(A)$ and a
homogeneous ideal $I\subseteq A$, generated in positive degrees, whose
projective zero set is $Y$. Its affine zero set includes the cone vertex.
The open affine-cone complement
$V=\operatorname{Spec}(A)\setminus V(I)$ is the principal
$\mathbb G_m$-bundle associated with $\mathcal O_U(-1)$.
For a finite graded $A$-module $M$ representing $E$, affineness of
$V\to U$ and the local-cohomology sequence give, for $p\geq1$,
\[
 H_I^{p+1}(M)
 \cong\bigoplus_{a\in\mathbb Z}H^p(U,E(a)|_U).
 \tag{C.12.1}
\]
We may choose $M$ finite, and replace it by $M\oplus A/A_+$ without
changing its associated sheaf or (C.12.1). Then $M$ has full support:
every minimal prime of the standard graded ring $A$ is homogeneous;
a minimal prime other than the vertex $A_+$ belongs to
$\operatorname{Proj}(A)$ and lies in the support of $M$ because $E$ has
full support, while $A/A_+$ supplies the vertex. A closed support containing
all minimal primes is all of $\operatorname{Spec}(A)$.
The empty $X$ case is immediate and may be omitted from this construction.

Let $c=\operatorname{cd}_I(A)$, the largest degree with
$H_I^c(A)\ne0$. A finite flat Cech complex on generators of $I$ computes
local cohomology. Tensoring a free resolution of an arbitrary module $N$
with that complex gives a spectral sequence with terms
$\operatorname{Tor}_a^A(N,H_I^b(A))$ and total degree $b-a$.
Consequently $H_I^j(N)=0$ for $j>c$ and
\[
 H_I^c(N)\cong N\otimes_A H_I^c(A).
 \tag{C.12.2}
\]
If $c\leq m$, (C.12.1), applied to finite graded modules representing
coherent extensions of sheaves on $U$, proves (1).
Such extensions exist because $X$ is noetherian and $U$ is quasi-compact;
one may construct one by extending finitely many generators and relations
on a finite affine cover and clearing denominators on the overlaps.

Suppose instead $c\geq m+1$. Put $Q=H_I^c(A)$.
By C.011, (C.12.2) gives $H_I^c(M)=M\otimes_A Q\ne0$.
Because $c>0$ and $H_I^c(A/I)=0$, the exact sequence of $0\to I\to A\to A/I\to0$
shows that $H_I^c(I)\to Q$ is onto. Under (C.12.2) this is multiplication
$I\otimes Q\to Q$. Hence $IQ=Q$, and also
$I H_I^c(M)=H_I^c(M)$.
A nonzero graded module $L$ satisfying $IL=L$, when $I$ is generated in
positive degrees, cannot be bounded below: a nonzero homogeneous element
of least degree could not be a sum of positive-degree elements times
homogeneous elements of smaller degree. Thus $H_I^c(M)$ has nonzero
components in arbitrarily negative degrees.
But (2) and (C.12.1), with $p=c-1\geq m$, say exactly that those components
vanish in all sufficiently negative degrees. This contradiction proves (1).
\end{proof}

\begin{saresult}{C.013}{Known}{Projective and formal duality in the form used here}
Let $X$ be a projective Cohen--Macaulay $k$-scheme, pure of dimension $r$,
let $\omega_X$ denote its dualizing sheaf, let $Y$ be closed, and let
$\widehat X$ be the completion along $Y$.
For every coherent $F$ and every integer $p$ there are natural identifications
\[
 H^p(X,F)^\vee\cong\operatorname{Ext}^{r-p}_X(F,\omega_X),
 \qquad
 H_Y^p(X,F)^\vee\cong
 \operatorname{Ext}^{r-p}_{\widehat X}(\widehat F,\widehat\omega_X).
 \tag{C.13.1}
\]
Here $(-)^\vee=\operatorname{Hom}_k(-,k)$ is the full algebraic dual,
and the formal Ext is hyper-Ext of coherent modules on the formal scheme.
The dual of $H_Y^p(X,F)\to H^p(X,F)$ is the completion map on the right.
Moreover
\[
 R\mathcal Hom_X(\omega_X,\omega_X)\simeq\mathcal O_X,
 \qquad
 R\mathcal Hom_{\widehat X}(\widehat\omega_X,\widehat\omega_X)
 \simeq\mathcal O_{\widehat X}.
 \tag{C.13.2}
\]
\end{saresult}
\begin{proof}[Derivation from the cited duality theorems]
The first formula and full support of $\omega_X$ are Stacks Project,
Lemma 48.27.5, tag 0FVZ. Biduality is Lemma 48.2.5, tag 0A89;
it gives (C.13.2), also after completion.
For completeness we spell out the projective reduction of the formal formula.
Embed $i:X\hookrightarrow P=\mathbb P_k^d$ and set $G=i_*F$.
Let $\mathcal J$ define $Y$ in $P$.
Local cohomology and duality on the regular projective space give
\[
 \begin{split}
 H_Y^p(P,G)&=\varinjlim_a
       \operatorname{Ext}^p_P(\mathcal O_P/\mathcal J^a,G),\\
 H_Y^p(P,G)^\vee&=\varprojlim_a
       \operatorname{Ext}^{d-p}_P
       (G,(\mathcal O_P/\mathcal J^a)\otimes\omega_P).
 \end{split}
\]
The second identification uses the Serre duality functor
$-\otimes\omega_P[d]$ on perfect complexes; both coherent arguments are
perfect on the regular scheme $P$. Duality on projective space is
SGA~2, XII, Theorem 1.1.
Resolve $G$ by a bounded complex of finite locally free sheaves on $P$.
On a finite affine cover with affine intersections, the complexes computing
the Ext groups above are inverse systems of Cech complexes. Their cohomology
groups are finite-dimensional over $k$ and therefore satisfy the
Mittag--Leffler condition. The inverse-limit spectral sequence has no
$\varprojlim^1$ contribution. Their inverse limit computes
$R\operatorname{Hom}_{\widehat P}(\widehat G,\widehat\omega_P)$.
Closed-immersion duality identifies
$i^!(\omega_P[d])$ with $\omega_X[r]$, since $X$ is Cohen--Macaulay and pure-dimensional;
the same identification holds after the flat completion.
This proves the formal formula and its compatibility with the completion
map. It is also the specialization to $\operatorname{Spec}(k)$ of
Alonso Tarrio--Jeremias Lopez--Lipman, \S5.3, formula (5.3.1), applied
to the dual of $F$. No invertibility of $\omega_X$ has been used.

The source locators used above are \cite[48.27.5 and 48.2.5]{C-stacks}; \cite[5.3, especially (5.3.1)]{C-ajl1997}.
\end{proof}

\begin{saresult}{C.014}{Proved}{The entire corrected chain on cote 125, scan 003}
Assume the hypotheses of C.013, choose a very ample $\mathcal O_X(1)$,
and fix $0<m<r$. Put $d=r-m>0$ and $U=X\setminus Y$.
The following seven conditions are equivalent:
\begin{enumerate}
\item $H^p(U,F)=0$ for every coherent $F$ on $U$ and every $p\geq m$.
\item $H^p(U,\mathcal O_X(-N)|_U)=0$ for $p\geq m$, $N\gg0$.
\item For every coherent $F$ on $X$, $H_Y^p(X,F)\to H^p(X,F)$ is
bijective for $p>m$ and surjective for $p=m$.
\item Condition (3) holds for $F=\mathcal O_X(-N)$, $N\gg0$.
\item For every coherent $F$ on $X$, the completion map
\[
 \operatorname{Ext}^q_X(F,\omega_X)\longrightarrow
 \operatorname{Ext}^q_{\widehat X}(\widehat F,\widehat\omega_X)
\]
is bijective for $q<d$ and injective for $q=d$.
\item For $N\gg0$, $H^0(X,\mathcal O_X(N))\to
H^0(\widehat X,\mathcal O_{\widehat X}(N))$ is bijective and
$H^q(\widehat X,\mathcal O_{\widehat X}(N))=0$ for $0<q<d$.
\item For every finite locally free $E$ on $X$, the completion map
$H^q(X,E)\to H^q(\widehat X,\widehat E)$ is bijective for $q<d$
and injective for $q=d$.
\end{enumerate}
The same assertions include $Y=\varnothing$ and $Y=X$; no nonemptiness
hypothesis has been suppressed. In (4), the surjectivity requirement at
$p=m$ is automatic for $N\gg0$.
\end{saresult}
\begin{proof}
C.012 applied to $E=\mathcal O_X$ proves (1)$\Leftrightarrow$(2).
For each $F$, the localization sequence
\[
 H_Y^p(X,F)\longrightarrow H^p(X,F)\longrightarrow H^p(U,F)
 \longrightarrow H_Y^{p+1}(X,F)\longrightarrow H^{p+1}(X,F)
\]
proves (1)$\Leftrightarrow$(3) and (2)$\Leftrightarrow$(4).
Extensions of coherent sheaves from $U$ to $X$ ensure that this is the same
quantification over coherent sheaves in (1) and (3).
C.013, together with the exactness and faithfulness of algebraic duality
of vector spaces, proves (3)$\Leftrightarrow$(5), including the boundary
surjection/injection and the precise index $q=r-p$.

For a locally free $E$, put $F=E^\vee\otimes\omega_X$.
By (C.13.2),
$R\mathcal Hom_X(F,\omega_X)=E$, and the same holds formally.
Thus (5) implies (7). Conversely, (7) implies (6) by taking $E=\mathcal O_X(N)$
and using Serre vanishing on $X$; the injectivity at $q=d>0$ is automatic
for those twists because its source is zero.

Finally assume (6). Serre vanishing makes it equivalent to the isomorphism
and injection assertions for the same positive twists through degree $d$.
Apply (C.13.1)--(C.13.2) with $F=\omega_X(-N)$.
We obtain the local-support comparison of (3) for these $F$, and therefore
\[
 H^p(U,\omega_X(-N)|_U)=0\qquad(p\geq m, N\gg0).
\]
The dualizing sheaf has support all of $X$; C.012 now proves (1).
This is the missing argument in the earlier study's purported
Cohen--Macaulay proof: $E^\vee\otimes\omega_X$ is not assumed locally free.

For the last assertion, Serre duality gives
\[
 H^m(X,\mathcal O_X(-N))^\vee
 =H^{r-m}(X,\omega_X(N))=0\qquad(N\gg0),
\]
because $r-m>0$. Thus every map into this zero group is surjective.
When $Y=X$, all completion maps and supported-to-ordinary maps are
identities and $U$ is empty. When $Y=\varnothing$ and $X$ is nonempty,
condition (6) fails because positive powers of a very ample line bundle
have nonzero sections while the completion is empty, and the equivalence
still applies. If $X$ is empty, every listed condition is vacuous.
\end{proof}

\begin{saresult}{C.015}{Refuted}{The positive-twist reading of the top test box}
Replacing $\mathcal O_X(-N)$ by $\mathcal O_X(N)$ in C.014(2) is false,
even for smooth projective $X$, nonempty $Y$, and $0<m<\dim X$.
Likewise $\Omega^r_{X/k}$ cannot be used as a synonym for $\omega_X$
for a general Cohen--Macaulay scheme.
\end{saresult}
\begin{proof}
For the first assertion take
$X=\mathbb P^1_k\times\mathbb P^1_k$, $Y=\mathbb P^1_k\times\{\infty\}$,
$\mathcal O_X(1)=\mathcal O(1,1)$, and $m=1$.
Then $U=\mathbb P^1\times\mathbb A^1$ and
\[
 H^p(U,\mathcal O_X(N)|_U)
 =H^p(\mathbb P^1,\mathcal O(N))\otimes k[t]=0
 \quad(p\geq1, N\geq0).
\]
But $H^1(U,\mathcal O(-2,0)|_U)=k[t]\ne0$, so
$\operatorname{cd}(U)=1$.
The formulas follow directly from the standard two-affine Cech complex
for $\mathbb P^1$; hence all hypotheses and the failed conclusion are explicit.

For the second assertion let $\operatorname{char}(k)\ne2$,
$A=k[\epsilon]/(\epsilon^2)$, and
$X=\mathbb P^2_k\times\operatorname{Spec}(A)$.
This is projective, Gorenstein and pure of dimension two.
Its dualizing sheaf is $\omega_{\mathbb P^2}\otimes_k A$, whereas
\[
 \Omega^2_{X/k}
 =\bigl(\Omega^2_{\mathbb P^2/k}\otimes_k A\bigr)
 \oplus\bigl(\Omega^1_{\mathbb P^2/k}\otimes_k k\,d\epsilon\bigr).
\]
The Euler sequence gives $H^1(\mathbb P^2,\Omega^1)=k$ and
$H^1(\mathbb P^2,\Omega^2)=0$.
Thus $\operatorname{Ext}^1_X(\mathcal O_X,\Omega^2_{X/k})=k$, but
$H^1(X,\mathcal O_X)^\vee=0$. The claimed Serre duality with
$\Omega^2_{X/k}$ fails. C.014 uses the actual dualizing sheaf;
for a smooth $X$ it agrees with $\Omega^r_{X/k}$.
\end{proof}

\begin{saresult}{C.016}{Proved}{The question mark on scan 006: formal Ext and inverse limits}
Let $X$ be projective over a field, let $\mathcal I$ be a coherent ideal,
and let $F,G$ be coherent $\mathcal O_X$-modules. Put
$F_a=F/\mathcal I^aF$ and let hats denote completion along
$V(\mathcal I)$. Then for every $q\geq0$ the canonical map is an
isomorphism
\[
 \operatorname{Ext}^q_{\widehat X}(\widehat G,\widehat F)
 \xrightarrow{\ \sim\ }
 \varprojlim_a\operatorname{Ext}^q_X(G,F_a).
 \tag{C.16.1}
\]
No Cohen--Macaulay or local-freeness hypothesis is needed for this
comparison itself. Such hypotheses enter only when these groups are
identified with dual cohomology groups.
\end{saresult}
\begin{proof}
A coherent sheaf on a projective scheme is the quotient of a finite
sum of negative powers of a very ample line bundle. Repeating this on
the coherent kernels constructs a resolution $E_\bullet\to G$ by
finite locally free sheaves. It may be infinite to the left; this
causes no convergence problem in the first-quadrant Hom complex.
Take a finite affine cover of $X$. Since $X$ is separated, all its
finite intersections are affine. The total Cech--Hom complex $C_a$
of $\mathcal Hom(E_\bullet,F_a)$ computes
$R\operatorname{Hom}_X(G,F_a)$.
The maps of the $F_a$ are onto. Exactness of sections of quasi-coherent
sheaves on affine schemes therefore makes the maps of $C_a$ onto
term by term. For a fixed total degree there are only finitely many
summands. The inverse limit of these complexes is the analogous
Cech--Hom complex on $\widehat X$: completion is flat and preserves
the resolution, and coherent modules on an affine noetherian formal
scheme have no higher cohomology.

For a countable inverse system of complexes the standard product
description of $R\varprojlim$ gives
\[
 0\to\varprojlim{}^1 H^{q-1}(C_a)\to
 H^q(R\varprojlim C_a)\to\varprojlim H^q(C_a)\to0.
\]
Termwise surjectivity identifies $R\varprojlim C_a$ with
$\varprojlim C_a$. Each $\operatorname{Ext}^q_X(G,F_a)$ is finite
dimensional over the field, by projectivity and coherence. The images
in any fixed term form a decreasing sequence of subspaces and
stabilize. Thus the cohomology inverse systems are Mittag--Leffler,
and their $\varprojlim{}^1$ vanishes. The displayed exact sequence
proves (C.16.1), naturally in both arguments. Applying this with the
primed scheme and $q=r-1-i$ answers the question mark on scan 006
whenever the primed scheme is projective. It proves more than the
earlier study's locally free version.
\end{proof}

\begin{saresult}{C.018}{Proved under stated hypotheses}{A complete Cartier-section induction with the boundary injection}
Let $X$ be projective Cohen--Macaulay and pure of dimension $r$.
Let $D_1,\ldots,D_c$ be effective ample Cartier divisors, with
$1\leq c<r$, whose local equations form a regular sequence, and let
$Y=D_1\cap\cdots\cap D_c$. For every locally free $E$, completion
along $Y$ gives an isomorphism
\[
 H^q(X,E)\longrightarrow H^q(\widehat X,\widehat E)
       \quad(q<r-c)
\]
and an injection for $q=r-c$.
The same result is obtained at every stage of the flag, with $c$
replaced by its length. The qualification ``ample'' cannot be dropped.
\end{saresult}
\begin{proof}
Write $L_j=\mathcal O_X(D_j)$ and let $s_j$ define $D_j$.
For $a\geq1$, the powered sequence $s_1^a,\ldots,s_c^a$ is still a
regular sequence. Its ideal $\mathcal J_a$ defines a thickening $Y_a$.
These ideals are cofinal with the powers of
$\mathcal I=(s_1,\ldots,s_c)$, because
\[
 \mathcal I^{c(a-1)+1}\subseteq\mathcal J_a\subseteq\mathcal I^a.
\]
The Koszul complex resolves $E|_{Y_a}$ with term in degree $-j$
equal to
\[
 \bigoplus_{|J|=j} E\otimes
       \left(\bigotimes_{\nu\in J}L_\nu\right)^{-a},
       \qquad 0\leq j\leq c.
\]
For each nonempty $J$, the product line bundle in parentheses is
ample. Serre duality and Serre vanishing imply that, for $a$ large,
each displayed negative-twist sheaf has zero cohomology in every
degree below $r$. There are only finitely many $J$, so one bound on
$a$ works for all of them.

The quotient of the Koszul resolution by its degree-zero subcomplex
$E$ represents the cone of $E\to E|_{Y_a}$ and has terms only in
degrees $-c,\ldots,-1$. Its hypercohomology spectral sequence has
terms of total degree $v-j$, where $H^v$ is the cohomology of the
degree-$j$ negative-twist sheaf. All those terms are zero when
$v-j<r-c$. Thus the cone has no hypercohomology below $r-c$.
Its exact triangle proves the claimed comparison through the boundary
degree for $Y_a$. C.016 and cofinality pass this to the completion.
For a flag of length $j$, the same argument uses only the first $j$
Koszul factors. Adding the next divisor adds one factor and lowers
the bound by one, from $r-j$ to $r-j-1$. This specifies the induction
step, its initial case $j=1$, and its final bound.

To disprove the assertion without ampleness, take
$X=\mathbb P^1\times\mathbb P^1$, $D=\mathbb P^1\times\{0\}$,
and $E=\mathcal O(0,-1)$. The divisor is effective Cartier and
$X$ is smooth of dimension two. Its formal neighborhood is
$\mathbb P^1\times\operatorname{Spf}k[[t]]$, on which $E$ is
trivial. Hence $H^0(X,E)=0$ while
$H^0(\widehat X,\widehat E)=k[[t]]$. The required degree-zero
isomorphism fails.

The source locators used above are \cite[XII, 2.1--2.3, pp.~114--115]{C-sga2}.
\end{proof}

\paragraph{The Cartier interpretation is an additional theorem.}
The handwritten page 006 does not state an induction conclusion and
does not define $X'$ or the subscripts $a$ and $-$. Its top row has
the form of localization at a point $a$, with $X'=X\setminus\{a\}$.
The shift $r-1$ then comes from punctured-local duality; see C.021.
It is \emph{not} the Cartier-section sequence. For a Cartier section
$D$, the actual sequence is
$0\to G(-D)\to G\to i_*G|_D\to0$, when $G$ is locally free.
It has an essential twist $G(-D)$ in the third term after applying
$R\mathcal Hom(F,-)$, and adjunction gives
$\omega_D=\omega_X(D)|_D$. Thus C.018 supplies a complete and useful
Cartier induction, but cannot be attributed as the literal meaning
of the diagram. C.021--C.022 prove the local reading and its nested
support induction instead.

More generally the exact obstruction to an attempted induction can
be stated without choosing a reading. Put
$C_k=\operatorname{Cone}(R\Gamma(X,E(kD))\to
R\Gamma(\widehat X,\widehat E(kD)))$ and define $D_k$ similarly for
$D$ and $E(kD)|_D$. The divisor sequence gives
$C_{k-1}\to C_k\to D_k\to C_{k-1}[1]$.
If $H^i(D_k)=0$ for $i<d-1$, all $k$, and $H^i(C_N)=0$ for
$i<d$ and $N$ sufficiently large, then $H^i(C_k)=0$ for $i<d$
and all $k$: the long exact sequence makes each transition
$H^i(C_{k-1})\to H^i(C_k)$ injective for $i<d$, and iteration
embeds it in a zero group. This is the complete diagram chase.
The eventual vanishing is a necessary extra input to this argument,
not something that the unlabeled diagram proves by itself.

\begin{saresult}{C.019}{Proved}{The threshold on scan 008 and the scope of the diagonal margin}
Let $P=\mathbb P^r_k$, let $Y\subset P$ be closed, and let
$0<n<r$. The following are equivalent:
\begin{enumerate}
\item $H^0(P,\mathcal O_P(a))\to
H^0(\widehat P,\mathcal O_{\widehat P}(a))$ is an isomorphism and
$H^q(\widehat P,\mathcal O_{\widehat P}(a))=0$ for $0<q<n$,
for all sufficiently large $a$;
\item the same assertions hold for every integer $a$;
\item $\operatorname{cd}(P\setminus Y)\leq r-n-1$.
\end{enumerate}
Replacing $\mathcal O_P(a)$ by $\omega_P(a)$ gives the same criterion.
Thus the written $r-n$ on scan 008 is the first vanishing degree,
not the usual cohomological dimension. The bound is attained by
$Y=\mathbb P^n\subset\mathbb P^r$.

For a general pure Cohen--Macaulay projective $X$ of dimension $r$
and $\dim Y\leq n$, the diagonal margin on scan 003 has the correct
comparison
\[
 H^q(X,E)\xrightarrow{\sim}H^q(X\setminus Y,E)
       \quad(q<r-n-1)
 \tag{C.19.1}
\]
for locally free $E$, including $q=0$ when $n<r-1$.
The further claim that these groups vanish for $0<q<r-n-1$
needs vanishing on $X$ itself and is false for general such $X$.
\end{saresult}
\begin{proof}
Apply C.014 with $m=r-n$. Condition (1) is its positive-twist test.
Its equivalent locally free comparison gives (2), since
$H^q(\mathbb P^r,\mathcal O(a))=0$ for $0<q<r$ and every $a$.
Conversely (2) includes (1), and the cohomological bound in C.014 is
$\operatorname{cd}<m$, which is $\operatorname{cd}\leq r-n-1$.
The formula $\omega_P=\mathcal O_P(-r-1)$ merely reindexes the twists.
For a linear $Y=\mathbb P^n$, the complement is covered by the
$r-n$ standard affine opens whose coordinates do not all vanish.
Their intersections are affine, so Cech cohomology gives the upper
bound $r-n-1$. A complementary linear $\mathbb P^{r-n-1}$ is a
closed subvariety of the complement. The direct image of
$\mathcal O_{\mathbb P^{r-n-1}}(-(r-n))$ has nonzero cohomology in
degree $r-n-1$, by projective-space duality. This proves equality.

At a point of $Y$, Cohen--Macaulayness and pure dimension give
depth at least $r-n$. The same is true for a locally free $E$.
The local-cohomology depth criterion gives
$\mathcal H^j_Y(E)=0$ for $j<r-n$. The local-to-global spectral
sequence then gives $H^j_Y(X,E)=0$ in that range. The localization
sequence proves (C.19.1).
For the failure of the asserted vanishing, take
$X=C\times\mathbb P^2$ with $C$ a smooth plane cubic and take
$Y$ to be one closed rational point, over an algebraically closed
field. Here $r=3$, $n=0$, and
$H^1(X,\mathcal O_X)=k$ by the cubic exact sequence on
$\mathbb P^2$ and the Kunneth formula. Equation (C.19.1) gives
$H^1(X\setminus Y,\mathcal O_X)=k\ne0$, in the stated vanishing
range $0<1<2$. This verifies the missing hypothesis rather than
silently restricting the margin to projective space.

The source locators used above are \cite[scans 003 and 008]{C-ms125}.
\end{proof}

\begin{saresult}{C.021}{Proved under stated hypotheses}{The punctured-local diagram of scan 006 and its complete induction}
Let $(A,\mathfrak m)$ be a complete noetherian local Cohen--Macaulay
ring of dimension $r$, with canonical module $\omega_A$.
Put $X=\operatorname{Spec}A$, $a=V(\mathfrak m)$,
$U=X\setminus\{a\}$, and let $I\subseteq\mathfrak m$,
$Y=V(I)$ and $Y'=Y\cap U$. Let $P$ be a bounded complex with
finite cohomology modules on $X$ and
put $Q=R\operatorname{Hom}_A(P,\omega_A)$.
For the injective hull $E_A$ of the residue field, write
$\mathbb D=R\operatorname{Hom}_A(-,E_A)$.
Completion with a hat below is along $Y$, not along a newly chosen
closed point. There are natural isomorphisms
\begin{align}
 \mathbb D P&\simeq R\Gamma_a(X,Q)[r],&
 \mathbb D R\Gamma_a(X,P)&\simeq Q[r],\tag{C.21.1}\\
 \mathbb D R\Gamma(U,P)&\simeq R\Gamma(U,Q)[r-1],\tag{C.21.2}\\
 \mathbb D R\Gamma_Y(X,P)&\simeq
 R\Gamma_a(\widehat X,\widehat Q)[r],\tag{C.21.3}\\
 \mathbb D R\Gamma_{Y'}(U,P)&\simeq
 R\Gamma(\widehat U,\widehat Q)[r-1].\tag{C.21.4}
\end{align}
They identify the dual of
\[
 R\Gamma_a(X,P)\longrightarrow R\Gamma_Y(X,P)
 \longrightarrow R\Gamma_{Y'}(U,P)\longrightarrow
 \tag{C.21.5}
\]
with the formal localization triangle, with its arrows reversed
and the displayed shifts. In particular, for perfect $F$ and
bounded coherent $G$, take
$P=R\mathcal Hom(F,G)$ and
$Q=R\mathcal Hom(G,F\otimes^L\omega_A)$.
When $A$ is Gorenstein and a trivialization $\omega_A\simeq A$ is
chosen, these are exactly the swapped-argument Ext groups drawn on
scan 006, with $X'=U$ and subscript $a$ denoting support at the
closed point. The $r-1$ is forced by localization.

More generally, for closed $Y\subseteq Z\subseteq X$ containing
$a$, the dual of the localization triangle
$R\Gamma_Y P\to R\Gamma_Z P\to
R\Gamma_{Z\setminus Y}(X\setminus Y,P)$ is the triangle determined
by the restriction
\[
 R\Gamma_a(\widehat X_Z,\widehat Q_Z)
 \longrightarrow R\Gamma_a(\widehat X_Y,\widehat Q_Y).
 \tag{C.21.6}
\]
Precisely, the dual of the relative term is
\[
 \mathbb D R\Gamma_{Z\setminus Y}(X\setminus Y,P)
 \simeq \operatorname{Cone}\bigl(
 R\Gamma_a(\widehat X_Z,\widehat Q_Z)\to
 R\Gamma_a(\widehat X_Y,\widehat Q_Y)\bigr)[r-1].
\]
Iteration proves the same assertion for every successive difference
of a finite closed filtration, including all endpoint connecting
maps. This is the complete localization induction visible in the
last two rows of the page.
\end{saresult}
\begin{proof}
Local duality, Stacks Project, Theorem 47.18.3 (0A84) and Lemma
47.18.4 (0AAK), gives
$\mathbb D R\Gamma_a(P)=R\operatorname{Hom}_A(P,\omega_A[r])$
because finite modules over $A$ are complete. The canonical
module is dualizing, so biduality gives
$R\operatorname{Hom}_A(Q,\omega_A)=P$; this is also
Stacks Project, Lemma 48.2.5 (0A89), with the dualizing
shift canceled twice. Apply local duality to $Q$ as well.
Its local cohomology modules are artinian, hence Matlis
reflexive because $A$ is complete. Dualizing the resulting
identity $\mathbb D R\Gamma_aQ=P[r]$ gives
\[
 \mathbb D P=R\Gamma_aQ[r].
\]
These identifications use the natural bidual evaluation maps,
so they are compatible with the localization evaluation map.
For perfect $P$ one can see this directly from
$\mathbb D P=P^\vee\otimes E_A$ and
$R\Gamma_a\omega_A[r]=E_A$.
This proves (C.21.1). Dualize
$R\Gamma_aP\to P\to R\Gamma(U,P)$ and use these
identifications. Its middle arrow becomes
$R\Gamma_aQ[r]\to Q[r]$, so its remaining term is
$R\Gamma(U,Q)[r-1]$. This proves (C.21.2) and its arrows,
not merely its cohomology dimensions.

For completeness, here is the formal-completion step behind
(C.21.3). Let $C_I$ be the finite flat alternating Cech complex
on generators of $I$, including $A$ in degree zero. Local
cohomology is $C_I\otimes_A^L-$ and derived completion is
$R\operatorname{Hom}_A(C_I,-)$; these are the explicit complexes
of Stacks Project, Lemma 47.12.3 (0A6Y), and its cited completion
formula. Therefore
\[
 \mathbb D R\Gamma_I P
 =R\operatorname{Hom}_A(C_I,\mathbb D P)
 =L\Lambda_I R\Gamma_a Q[r].
 \tag{C.21.7}
\]
The last unshifted term is precisely
$R\Gamma_a(\widehat X,\widehat Q)$.
Here is a direct verification of that identification, which also
specifies the support convention. A finite affine cover of
$U$ by principal opens and its augmented Cech complex compute
cohomology with support at $a$. Apply this complex to
$Q\otimes_A^L A/I^b$ and take $R\varprojlim_b$.
This is $L\Lambda_I R\Gamma_aQ$ because the Cech terms are flat
and there are finitely many of them. A bounded-above resolution of
$Q$ by finite free modules has, in each fixed degree, a pro-zero
negative Tor obstruction to replacing the derived quotient by the
termwise coherent quotient. This is Artin--Rees applied to the
kernels and images in a finite initial segment of the resolution:
for any fixed power $I^b$, a sufficiently large power $I^c$ sends
each kernel inside the required $I^b$-multiple, so the induced
transition on the relevant Tor group is zero. The Cech complex
has finite length, hence only finitely many such degrees affect
any chosen cohomology degree. Its derived inverse limit is
unchanged by these pro-zero systems. On each affine formal open,
ordinary completion of a coherent module is exact and equals
its inverse limit of quotients. The resulting augmented complex
is exactly the formal Cech complex with support at $a$.
This proves the identification and (C.21.3).

Any finite $A$-module is also $I$-adically complete. Indeed,
an $I$-adic Cauchy sequence is $\mathfrak m$-adic Cauchy since
$I^b\subseteq\mathfrak m^b$; its limit remains in each specified
coset modulo $I^b$ because $I^bM$ is $\mathfrak m$-adically
closed, by Krull intersection for the finite module $M/I^bM$.
Thus $R\Gamma(\widehat X,\widehat Q)=Q$.
Dualizing (C.21.5), and using (C.21.1) and (C.21.3), now gives
\[
 \mathbb D R\Gamma_{Y'}(U,P)\longrightarrow
 R\Gamma_a(\widehat X,\widehat Q)[r]\longrightarrow Q[r]
 \longrightarrow.
\]
The formal support triangle identifies its first term with
$R\Gamma(\widehat U,\widehat Q)[r-1]$, proving (C.21.4).
C.022 below identifies its cohomology with the inverse limits
drawn at the question mark whenever the arguments there are
coherent sheaves. The same argument for $Y\subseteq Z$ proves
(C.21.6). For a finite flag, start with one closed inclusion and
add one more. Localization supplies an exact triangle and the
dual restriction supplies the corresponding exact triangle on
formal completions. The first two comparison maps are the
already constructed isomorphisms (C.21.3), so their induced map
on cones is an isomorphism. This proves the next relative piece
and all its connecting maps. The empty relative piece has zero
complex and is the initial case. This is an actual induction,
with no unstated vanishing or positivity hypothesis.

The source locators used above are \cite[47.18.3--47.18.4, 48.2.5, 47.12.3]{C-stacks}; \cite[XIII, Section~1 after Problem~1.4, p.~139]{C-sga2}; \cite[1.2]{C-szpiro1976}.
\end{proof}

\begin{saresult}{C.022}{Proved}{The question mark also holds on a punctured complete local scheme}
Let $(A,\mathfrak m)$ be a complete noetherian local ring, put
$U=\operatorname{Spec}A\setminus\{\mathfrak m\}$, and let
$I\subseteq\mathfrak m$. Let $F,G$ be finite $A$-modules and
write $F_b=(F/I^bF)|_U$. If $\widehat U$ is the formal
completion along $V(I)\cap U$, then, for every $q\geq0$,
\[
 \operatorname{Ext}^q_{\widehat U}(\widehat G,\widehat F)
 \xrightarrow{\ \sim\ }
 \varprojlim_b\operatorname{Ext}^q_U(G|_U,F_b).
 \tag{C.22.1}
\]
No regularity, Cohen--Macaulayness or local freeness is required
for this inverse-limit statement.
\end{saresult}
\begin{proof}
Use a resolution of $G$ by finite free $A$-modules and a finite
principal-affine cover of $U$. As in C.016, its Cech--Hom
complexes $C_b$ have termwise surjective transition maps, and
their inverse limit computes formal Ext. It remains to prove
$\varprojlim{}^1H^{q-1}(C_b)=0$.

Each $N_b=\operatorname{Ext}^j_U(G|_U,F_b)$ is Matlis reflexive.
To see this, the localization exact sequence writes it as an
extension of a quotient of the finite module
$\operatorname{Ext}^j_A(G,F/I^bF)$ by a submodule of
$\operatorname{Ext}^{j+1}_{\mathfrak m}(G,F/I^bF)$.
The latter module is artinian: the spectral sequence with terms
$H^v_{\mathfrak m}(\operatorname{Ext}^w_A(G,F/I^bF))$ has
artinian terms and only finitely many terms in any total degree.
The needed artinianness of local cohomology of finite modules
also follows at once from local duality over the complete local
ring, which has a dualizing complex, and Matlis duality.
Finite and artinian modules are Matlis reflexive by Stacks
Project, Proposition 47.7.8 (08Z9). Exactness of Matlis duality
and the five lemma show reflexivity is preserved by extensions.
Thus $N_b\simeq DDN_b$ canonically.

For any inverse sequence of Matlis-reflexive modules $N_b$, the
map $1-\mathrm{shift}:\prod_bN_b\to\prod_bN_b$ is surjective.
Indeed it is the Matlis dual of the injective telescope map
$1-\mathrm{shift}:\bigoplus_bDN_b\to\bigoplus_bDN_b$;
injectivity follows by looking at the first nonzero component
of a finite-support vector. Matlis duality is exact and carries
the direct sums to the products of $DDN_b=N_b$.
The product description of $\varprojlim{}^1$ therefore gives
$\varprojlim{}^1N_b=0$. Applying this to the cohomology of $C_b$
and then the Milnor exact sequence proves (C.22.1).
This argument explains why projectivity in C.016 can be replaced
here by completeness at the puncture; neither hypothesis may
simply be discarded in a general inverse-limit argument.

The source locators used above are \cite[47.7.8]{C-stacks}; \cite[XIII, Section~1 after Problem~1.4, p.~139]{C-sga2}.
\end{proof}

\paragraph{The literal scope and the two readings.}
For $X'=X\setminus\{a\}$, C.021 identifies the top and bottom
rows and C.022 resolves the question mark. For a projective
Cartier section, C.016 and C.018 are the corrected theorem, with
its necessary twist and ampleness. These readings are not
interchangeable. For example, let $X=\mathbb P^2$, $X'$ a line,
$a$ a closed point, and $F=G=\mathcal O_X$. The alleged
localization row would contain
\[
 0=H^1(X,\mathcal O_X)\longrightarrow
 0=H^1(X',\mathcal O_{X'})\longrightarrow
 H^2_a(X,\mathcal O_X)\longrightarrow
 0=H^2(X,\mathcal O_X).
\]
Its middle term is nonzero: in regular local parameters $x,y$
at $a$, the class of $x^{-1}y^{-1}$ is nonzero in
$A_{xy}/(A_x+A_y)$. Hence this row is not exact. Also,
Cohen--Macaulayness alone does not permit swapping $F,G$ without
the canonical-module factor in C.021. The complete two-dimensional
ring $A=k[[s^3,s^2t,st^2,t^3]]$ is free with basis
$1,s^2t,st^2$ over $k[[s^3,t^3]]$, so it is Cohen--Macaulay.
Modulo $(s^3,t^3)$ it has a two-dimensional socle spanned by
$s^2t,st^2$ and radical square zero. Its canonical module
therefore needs two generators, whereas $A$ needs one.
Local duality gives $D H^2_{\mathfrak m}(A)=\omega_A\not\simeq A$.
Perfection cannot simply be removed from the swapped-argument
specialization, even when the ring is Gorenstein. Take
$A=k[[x,y,\epsilon]]/(\epsilon^2)$, which is a complete
two-dimensional Gorenstein local ring, and take $F=G=k$.
The tensor product of the Koszul resolution on $x,y$ with
the resolution
$\cdots\xrightarrow{\epsilon} k[\epsilon]/(\epsilon^2)
\xrightarrow{\epsilon} k[\epsilon]/(\epsilon^2)\to k$
is a free resolution of $k$ over $A$.
After applying $\operatorname{Hom}_A(-,k)$ all its
differentials are zero. Thus
$\dim_k\operatorname{Ext}^0_A(k,k)=1$ and
$\dim_k\operatorname{Ext}^2_A(k,k)=4$.
All these Ext groups are supported at $\mathfrak m$.
The untwisted, unrestricted swap on the top row of scan 006
would identify the dual of the first vector space with the
second one, which is impossible. Here
$R\operatorname{Hom}_A(k,k)$ is not a bounded coherent
complex, exactly the hypothesis excluded above.
Thus the untwisted swap fails even in dimension two. The theorem
above gives the exact correction and completes the formal
localization induction in its full Cohen--Macaulay scope.

\begin{saresult}{C.020}{Known}{The verified literature comparison and the endpoint injection}
For the seven conditions in C.014, the following attributions follow
from the opened primary texts. They distinguish the ambient scheme
from the closed subscheme and do not replace a source hypothesis by
Cohen--Macaulayness without proof.
\begin{enumerate}
\item SGA~2, XII, 1.1--1.4, pp.~109--111, supplies projective
Serre duality and the equivalence of eventual negative-twist
vanishing with depth at closed points. It is not the full
complement/formal-completion equivalence for arbitrary $Y$.
XII, Theorem~2.1, Corollary~2.2 and Lemma~2.3,
pp.~114--115 (original margins 143--144), give comparison for
an ample effective section regular on the coherent sheaf, in a
projective family, with the stated fibrewise depth and flatness
hypotheses. Over a field, depth $>d$ at closed points gives
isomorphisms below $d$ and a monomorphism at $d$. In particular,
for a pure Cohen--Macaulay $r$-fold, an ample Cartier divisor,
and a vector bundle, the bound is $d=r-1$. This is a sufficient
geometric hypothesis, not a converse for an arbitrary closed $Y$.
The proof and its precise scope are C.018 with $c=1$.

\item Hartshorne, \emph{Ample subvarieties of algebraic varieties}
(1970), III, Theorem~3.4(b), p.~96, with proof on pp.~97--98,
contains both directions of C.014(1)$\Longleftrightarrow$(7)
when $X$ is a complete nonsingular variety. If $\dim X=r$ and
the cohomological dimension bound is $\operatorname{cd}U\leq a$,
its comparison is an isomorphism for $q<r-a-1$ and injective
for $q=r-a-1$. Put $a=m-1$ to obtain exactly $q<r-m$ and
the boundary injection in C.014. III, Proposition~3.1,
p.~92, supplies the negative-twist test in the smooth
quasiprojective case; III, Theorem~3.3, pp.~94--96, supplies
formal duality. These statements use a nonsingular ambient
variety. C.011--C.014 give the argument for a projective
Cohen--Macaulay ambient scheme whose canonical module need
not be invertible.

\item Hartshorne, \emph{Cohomology of non-complete algebraic
varieties} (1971), Proposition~2.9, p.~261, assumes
$X=\mathbb P^r_k$ and condition $(I_s^*)$: for $i<s$,
$\bigoplus_{v\in\mathbb Z}H^i(\widehat X,
\mathcal O_{\widehat X}(v))$ is a finite free graded
$k[x_0,\ldots,x_r]$-module generated in degree zero.
For an integer $d\leq s$, the proposition gives both directions
between $\operatorname{cd}U<r-d$, vanishing in those degrees
for the single untwisted sheaf $\omega_X$, and the
isomorphisms
$H^q(X,\mathcal O_X)\to H^q(\widehat X,
\mathcal O_{\widehat X})$ for $q<d$.
It does not state an additional boundary injection. For the
range $0<d<r$ considered here that injection for
$\mathcal O_X$ is automatic, since $H^d(\mathbb P^r,
\mathcal O)=0$; it is not a substitute for the vector-bundle
boundary condition of C.014(7).
Theorem~2.1, p.~259, proves $(I_s^*)$ when $Y$ is purely
$s$-dimensional and smooth over a characteristic-zero field,
or Cohen--Macaulay over a field of positive characteristic.
Here Cohen--Macaulayness is a condition on $Y$, while the
ambient scheme remains projective space.
Proposition~2.2, p.~260, separately contains an injection
at its boundary degree $s-\operatorname{hd}(F)$ for the
comparison map constructed there. Proposition~1.1, p.~258,
gives a different equivalent criterion for the next-to-top
vanishing on a smooth projective ambient variety: the
Lefschetz condition and intersection with every divisor.

\item Speiser, \emph{Cohomological Dimension of Non-Complete
Hypersurfaces}, Inventiones~21 (1973), \S2, Lemma, p.~147,
gives both directions of the positive-twist test. In his
notation, $V$ is a smooth projective $r$-dimensional variety
over an algebraically closed field and $W$ is a nonempty
connected closed subset. The bound
$\operatorname{cd}(V\setminus W)\leq r-k$ is equivalent to
the maps
\[
 H^q(V,\mathcal O_V(v))\longrightarrow
 H^q(\widehat V,\mathcal O_{\widehat V}(v))
\]
being injective for $0\leq q\leq k-1$ and isomorphisms
for $0\leq q\leq k-2$, for $v\gg0$.
Put $k=r-m+1=d+1$: the boundary is exactly $q=d$,
and Serre vanishing turns the remaining conditions into
C.014(6). Thus the injection is printed explicitly in
the lemma, even though it becomes automatic for sufficiently
positive twists in positive degree.
Theorem~C, pp.~147--148, gives the descent by one in the
cohomological bound under a smooth hyperplane section,
assuming $(I_k^*)$ with $k\geq2$; this is more specialized
than the general equivalence in C.014.
\end{enumerate}
These attributions concern the sources actually opened. The exact
clauses of Hartshorne (1968) and Ogus (1973) are the separate
source-verification gate recorded in the unchecked file; neither
paper is declared read on the strength of a later citation.
\end{saresult}
\begin{proof}[Derivation and audit of the comparison]
For SGA~2, XII, 1.1 is Ext--Serre duality on projective space.
Corollary~1.3 and Corollary~1.4 convert eventual vanishing of
$H^i(F(-N))$ to depth at closed points; the paragraph following
1.4 explicitly extends this formulation to an arbitrary
projective scheme by a projective embedding. For an ample
section $t$ regular on $F$, the sequence
\[
 0\longrightarrow F(-N)\xrightarrow{t^N}F
 \longrightarrow F/t^NF\longrightarrow0
\]
then gives comparison with each sufficiently thick infinitesimal
neighborhood, including injection in the endpoint degree.
XII, 2.3 and the inverse-limit argument in 2.2 pass to formal
cohomology. This proves the stated specialization directly from
the opened result. The printed wording of XII, 2.1 says depth
at the points of the fibres; the closed-point depth condition
used here is the one explicitly sufficient in XII, 1.4--1.5
and in that proof. No depth assertion at generic points is
being imposed.

For the nonsingular complete ambient variety in Hartshorne
(1970), formal duality transposes comparison into the support
map. The localization exact sequence identifies bijectivity
above $m$ and surjectivity at $m$ of that support map with
vanishing on $U$ from degree $m$ onward. Dualizing the
endpoint surjection is precisely the injection in his
Theorem~3.4(b). This is the same exact-sequence argument as
C.014, with an invertible canonical module. Proposition~3.1
then reduces the test to negative twists. This also explains
why replacing a negative test by a positive test on the
\emph{open complement} is incorrect; C.015 is the explicit
counterexample.

In Hartshorne (1971), $(I_s^*)$ is what permits an untwisted
test. A free graded module generated in degree zero is
determined by its degree-zero part; the vanishing of that
part is equivalent to the vanishing of all twists. His
Proposition~2.9 applies formal duality to this observation.
The source's $r$ in that proposition has been renamed $d$
above to avoid confusing it with ambient dimension. The
inequality is $d\leq s$ in the original page image; OCR of
the same line loses the equality. His single test sheaf is
the canonical line bundle on projective space, not
$\Omega_X^r$ on an arbitrary singular ambient scheme.
The distinction is essential by C.015.

For Speiser's lemma, substituting $k=d+1$ into its two
printed index ranges proves the asserted correspondence.
Its proof reduces arbitrary coherent sheaves by a surjection
from negative twists tensored with the canonical line bundle,
and explicitly cites Hartshorne (1968), Proposition~4.3.
It cannot justify replacing that line bundle by a locally
free sheaf when the ambient scheme is merely
Cohen--Macaulay; the full-support argument C.012 is the
additional input used here. Two indexing slips in the
printed proof do not change the lemma: its discussion
labelled $k=0$ uses the $k=1$ Lichtenbaum case, and the
two applications in the proof of Theorem~B say $k=1$
although their bounds require $k=2$. Our substitutions
use the displayed lemma and the direct proof C.014,
not these slips.

Finally, SGA~2, XIII, \S1, pp.~135--139, explicitly poses
affine duality and finiteness problems. Conjectures~1.1--1.2
cannot be invoked as proofs of the general chain; the
edition's notes record that those literal formulations
are false. Conjecture~1.3 is a finiteness assertion for
complements of local complete intersections in projective
space, not the vanishing/comparison equivalence of C.014.
Problem~1.4 asks for commutation with inverse limits.
Immediately afterward, p.~139, the text records Auslander's
theorem: for a finite module $M$ over a complete local ring,
$\operatorname{Ext}^i_A(-,M)$ converts direct limits to
inverse limits. C.022 gives the needed punctured-scheme
inverse-limit result by an explicit Matlis-duality argument.
The notation $X'=X\setminus\{a\}$ appears at the beginning
of XIII, \S2 on the same page, consistent with C.021's
reading of the handwritten diagram.

The source locators used above are \cite{C-sga2,C-hartshorne1970,C-hartshorne1971,C-speiser1973}.
\end{proof}
